\documentclass[11pt]{article}
\usepackage{amsmath,amssymb,longtable}
\emergencystretch=4em
\begin{document}
\section*{P11: Current complete mathematical model (as of 2026-10-20)}
\section*{Basis and scope of the revision}According to the definitions on pages 4–7 of the source text, where each task is performed once per vehicle, and all tasks are carried out together, the condition i=j should be excluded. Additionally, the equation relating access to vehicles should be added, along with a prohibition on self-looping. The constraints related to energy consumption, payload capacity, time constraints, and queueing should also be retained. This modification is intended to clarify the business logic; it does not aim to maintain exact equivalence to the original text.




\section{Formulations on Page 5}

\subsection{Objective Function}
\begin{equation}
    \label{eq:1}
    \min \lambda_1 \sum_{(i,j) \in \mathbb{A}} d_{ij}x_{ij} + \lambda_2 \sum_{i \in \mathbb{O} \cup \mathbb{F}} (\tau_i^* - \tau_i). \tag{1}
\end{equation}

\subsection{Constraints}
\begin{equation}
    \label{eq:2}
    \sum_{j \in \mathbb{O} \cup \mathbb{F} \cup \{2n+1\}} x_{ij} = 1, \forall i \in \mathbb{O}, \tag{2}
\end{equation}

\begin{equation}
    \label{eq:3}
    \sum_{i \in \{0\} \cup \mathbb{O} \cup \mathbb{F}} x_{ij} = 1, \forall j \in \mathbb{O}. \tag{3}
\end{equation}

\begin{equation}
    \label{eq:4}
    \sum_{j \in \mathbb{O} \cup \mathbb{F} \cup \{2n+1\}} x_{ij} \le 1, \forall i \in \mathbb{F}. \tag{4}
\end{equation}

\newpage

\section{Formulations on Page 6}

\subsection{Constraints}
\begin{equation}
    \label{eq:5}
    \sum_{i \in \{0\} \cup \mathbb{O} \cup \mathbb{F}} x_{ij} = \sum_{i \in \mathbb{O} \cup \mathbb{F} \cup \{2n+1\}} x_{ji}, \forall j \in \mathbb{F}, \tag{5}
\end{equation}

\begin{equation}
    \label{eq:6}
    u_i + q_i - Q^L(1 - x_{ij}) \le u_j, \forall i \in \{0\} \cup \mathbb{O} \cup \mathbb{F}, j \in \mathbb{O} \cup \mathbb{F} \cup \{2n+1\}, \tag{6}
\end{equation}

\begin{equation}
    \label{eq:7}
    u_i \le Q^L, \forall i \in \mathbb{O} \cup \mathbb{F}. \tag{7}
\end{equation}

\begin{equation}
    \label{eq:8}
    z_{ik} \le z_{0k}, \forall i \in \mathbb{O}, k \in \mathbb{K}, \tag{8}
\end{equation}

\begin{equation}
    \label{eq:9}
    z_{ik} + z_{jk} \le 3 - x_{0i} - x_{0j}, \forall i, j \in \mathbb{O} \cup \mathbb{F}, i\ne j, k \in \mathbb{K}, \tag{9}
\end{equation}

\begin{equation}
    \label{eq:10}
    2x_{ij} \le \sum_{k \in \mathbb{K}} (z_{ik} + z_{jk}), \forall i \in \{0\} \cup \mathbb{O} \cup \mathbb{F}, j \in \mathbb{O} \cup \mathbb{F} \cup \{2n + 1\}, \tag{10}
\end{equation}

\begin{equation}
    \label{eq:11}
    z_{ik} + (x_{ij} - 1) \le z_{jk}, \forall i, j \in \mathbb{O} \cup \mathbb{F}, k \in \mathbb{K}. \tag{11}
\end{equation}

\begin{equation}
    \label{eq:12}
    \sum_{k \in \mathbb{K}} z_{ik} \le 1, \forall i \in \mathbb{O} \cup \mathbb{F}, \tag{12}
\end{equation}

\begin{equation}
    \label{eq:13}
    z_{ik} = z_{(i+n)k}, \forall i \in \mathbb{P}, k \in \mathbb{K}. \tag{13}
\end{equation}

\begin{equation}
    \label{eq:14}
    y_j \le y_i - \gamma_{ij} - M(x_{ij} - 1), \forall i \in \{0\} \cup \mathbb{O}, j \in \mathbb{O} \cup \mathbb{F} \cup \{2n + 1\}, \tag{14}
\end{equation}

\begin{equation}
    \label{eq:15}
    y_j \le Q^E - \gamma_{ij} - M(x_{ij} - 1), \forall i \in \{0\} \cup \mathbb{F}, j \in \mathbb{O} \cup \mathbb{F} \cup \{2n + 1\}, \tag{15}
\end{equation}

\begin{equation}
    \label{eq:16}
    y_i \le Q^E, \forall i \in \{0, 2n + 1\} \cup \mathbb{O} \cup \mathbb{F}. \tag{16}
\end{equation}

\begin{equation}
    \label{eq:17}
    \tau_i^* + s_i \le \tau_{i+n}, \forall i \in \mathbb{P}, \tag{17}
\end{equation}

\begin{equation}
    \label{eq:18}
    \tau_i \le \tau_i^*, \forall i \in \mathbb{O} \cup \mathbb{F}, \tag{18}
\end{equation}

\begin{equation}
    \label{eq:19}
    b_i \le \tau_i^* \le e_i, \forall i \in \{0, 2n + 1\} \cup \mathbb{O} \cup \mathbb{F}. \tag{19}
\end{equation}

\begin{equation}
    \label{eq:20}
    \tau_i^* + s_i + t_{ij} \le \tau_j + M(1 - x_{ij}), \forall i \in \{0\} \cup \mathbb{O}, j \in \mathbb{O} \cup \mathbb{F} \cup \{2n + 1\}, \tag{20}
\end{equation}

\begin{equation}
    \label{eq:21}
    \tau_j \le \tau_i^* + s_i + t_{ij} + M(1 - x_{ij}), \forall i \in \{0\} \cup \mathbb{O}, j \in \mathbb{O} \cup \mathbb{F}, \tag{21}
\end{equation}

\begin{equation}
    \label{eq:22}
    \tau_i^* + (Q^E - y_i)g + t_{ij} \le \tau_j + M(1 - x_{ij}), \forall i \in \mathbb{F}, j \in \mathbb{O} \cup \mathbb{F} \cup \{2n + 1\}, \tag{22}
\end{equation}

\begin{equation}
    \label{eq:23}
    \tau_j \le \tau_i^* + (Q^E - y_i)g + t_{ij} + M(1 - x_{ij}), \forall i \in \mathbb{F}, j \in \mathbb{O} \cup \mathbb{F}. \tag{23}
\end{equation}

\newpage

\section{Formulations on Page 7}

\subsection{Constraints}
\begin{equation}
    \label{eq:24}
    \tau_j - \tau_i \le M w_{ij}, \forall i, j \in \{0, 2n + 1\} \cup \mathbb{O} \cup \mathbb{F}, \tag{24}
\end{equation}

\begin{equation}
    \label{eq:25}
    w_{ij} + w_{ji} = 1, \forall i, j \in \{0, 2n + 1\} \cup \mathbb{O} \cup \mathbb{F}, i \neq j, \tag{25}
\end{equation}

\begin{equation}
    \label{eq:26}
    M(w_{ij} - 1) \le \tau_j^* - \tau_i^* - s_i, \forall i, j \in \mathbb{S}^C(c), i \neq j, c \in \mathbb{C}, \tag{26}
\end{equation}

\begin{equation}
    \label{eq:27}
    M(w_{ij} - 1) \le \tau_j^* - \tau_i^* - (Q^E - y_i)g, \forall i, j \in \mathbb{S}^F(f), i \neq j, f \in \mathbb{F}^0. \tag{27}
\end{equation}

\begin{equation}
    \label{eq:28}
    x_{ij}, w_{ij}, z_{ik} \in \{0, 1\}, \forall i \in \{0\} \cup \mathbb{O} \cup \mathbb{F}, j \in \mathbb{O} \cup \mathbb{F} \cup \{2n + 1\}, k \in \mathbb{K}. \tag{28}
\end{equation}

\begin{equation}
    \label{eq:29}
    y_i, \tau_i, \tau_i^*, u_i \ge 0, \forall i \in \{0, 2n + 1\} \cup \mathbb{O} \cup \mathbb{F}, \tag{29}
\end{equation}



\section*{Collections, belonging, and complete domains}
$N=\{0,2n+1\}\cup O\cup F$, $I=\{0\}\cup O\cup F$, $J=O\cup F\cup\{2n+1\}$. $P$ refers to the pickup operations, while the delivery pairing is $i+n$. $K$ represents the vehicles used for transportation. The example includes all physical service groups and charging stations located at the same location.
\[
x_{ij}\in\{0,1\}\ (i\in I,j\in J),\quad
w_{ij}\in\{0,1\}\ (i,j\in N),\quad z_{ik}\in\{0,1\}\ (i\in N,k\in K).
\]
\[
\sum_kz_{ik}=\sum_{j\in J}x_{ij}\quad(i\in O\cup F),\qquad x_{ii}=0.
\]
Formula 9 is applicable only for different instances of $i,j$. Parameters such as service duration, distance, energy consumption during travel, charging rate, limited capacity, time constraints, and common physical locations all serve as instance parameters. All non-negative continuous variables remain within the aforementioned range; $0\le y_i\le Q^E$.

\end{document}
